\begin{mainquestion}{superposition}{Case Study 01 - Superposition}

\questionpoints
Consider the circuit below.

\begin{center}
    \begin{circuitikz}

    % 1. Leftmost Branch: Resistor (1 Ohm) and Current ix
    \draw (0,3) to[R, l_=$1\,\Omega$, i=$i_x$] (0,0);

    % 2. Middle Branch 1: 1A Current Source
    \draw (3,0) to[I, l_=$1\,\text{A}$] (3,3);

    % 3. Middle Branch 2: 2A Current Source
    \draw (6,0) to[I, l=$2\,\text{A}$] (6,3);

    % 4. Rightmost Branch: 10 Ohm Resistor
    \draw (9,3) to[R, l_=$10\,\Omega$, v^=$v_x$] (9,0);

    % 5. Top Connections
    \draw (0,3) -- (3,3); 
    \draw (3,3) -- (6,3); 
    
    % 3V Voltage Source 
    \draw (6,3) to[V, l=$3\,\text{V}$] (9,3);

    % 6. Bottom Connections
    \draw (0,0) -- (3,0); 
    
    % Dependent Voltage Source (4ix) 
    \draw (6,0) to[cV, l=$4i_x$] (3,0);

    % Wire from 2A Source to 10 Ohm Resistor
    \draw (6,0) -- (9,0);

    \end{circuitikz}
\end{center}

% =========================================================
\begin{partquestion}{calculatevx}
{Calculate $v_x$ in the circuit using \textbf{superposition}.

\textit{Solutions using nodal analysis or source transformation as the primary method will receive zero points.}}

\begin{answer}[5 in]
        To find $v_x$ using superposition, we calculate the contribution of each independent source individually while turning the others off. The dependent source $4i_x$ remains active in all steps.
        
        \[ v_x = v_x^{(1)} + v_x^{(2)} + v_x^{(3)} \]
        
        \textbf{Case 1: $1\,\text{A}$ Source Active} \\
        (Turn OFF $2\,\text{A}$ source $\to$ Open; Turn OFF $3\,\text{V}$ source $\to$ Short)
        
        Current leaving down the $1\,\Omega$ resistor is $i_x$. Current entering from the source is $1\,\text{A}$. The remaining current $(1 - i_x)$ flows right, through the shorted top wire, down the $10\,\Omega$ resistor, and returns via the dependent source.
        
        Apply KVL to the outer loop (clockwise):
        \[
        \underbrace{1 \cdot i_x}_{\text{Voltage across } 1\Omega} + 0 - \underbrace{10(1-i_x)}_{\text{Voltage } v_x^{(1)}} - \underbrace{4i_x}_{\text{Dependent Source}} = 0
        \]
        \[
        i_x - 10 + 10i_x - 4i_x = 0 \quad \Rightarrow \quad 7i_x = 10 \quad \Rightarrow \quad i_x = \frac{10}{7}\,\text{A}.
        \]
        Calculating $v_x^{(1)}$:
        \[
        v_x^{(1)} = 10(1 - i_x) = 10\left(1 - \frac{10}{7}\right) = -\frac{30}{7}\,\text{V}.
        \]
        
        \medskip
        \textbf{Case 2: $2\,\text{A}$ Source Active} \\
        (Turn OFF $1\,\text{A}$ source $\to$ Open; Turn OFF $3\,\text{V}$ source $\to$ Short)
        
        Apply KCL at the top node (supernode): Current leaving down $1\,\Omega$ is $i_x$. Input current is $2\,\text{A}$. Current down $10\,\Omega$ is $i_{10} = 2 - i_x$.
        
        Apply KVL to the outer loop:
        \[
        1 \cdot i_x - 10(2 - i_x) - 4i_x = 0
        \]
        \[
        i_x - 20 + 10i_x - 4i_x = 0 \quad \Rightarrow \quad 7i_x = 20 \quad \Rightarrow \quad i_x = \frac{20}{7}\,\text{A}.
        \]
        Calculating $v_x^{(2)}$:
        \[
        v_x^{(2)} = 10(2 - i_x) = 10\left(2 - \frac{20}{7}\right) = -\frac{60}{7}\,\text{V}.
        \]
        
        \medskip
        \textbf{Case 3: $3\,\text{V}$ Source Active} \\
        (Turn OFF $1\,\text{A}$ and $2\,\text{A}$ sources $\to$ Open)
        
        This forms a single loop. KVL tracing the direction of $i_x$:
        \[
        -1i_x + 4i_x - 10i_x + 3 = 0
        \]
        \[
        -7i_x = -3 \quad \Rightarrow \quad i_x = \frac{3}{7}\,\text{A}.
        \]
        Calculating $v_x^{(3)}$:
        \[
        v_x^{(3)} = -10 i_x = -10\left(\frac{3}{7}\right) = -\frac{30}{7}\,\text{V}.
        \]
        
        \textbf{Total Voltage $v_x$}
        \[
        v_x = \left(-\frac{30}{7}\right) + \left(-\frac{60}{7}\right) + \left(-\frac{30}{7}\right) = -\frac{120}{7}\,\text{V} \approx -17.14\,\text{V}.
        \]
\end{answer}

\qnewpage
\end{partquestion}

% =========================================================
\begin{partquestion}{comparenodal}
{Compare your result from part (a) with the answer obtained using \textbf{nodal analysis} (refer to Homework 4, Question 4). Which method do you prefer?}

\begin{answer}
        We previously solved this problem in Homework 4, Question 4 using Nodal Analysis. Using that method, we found the node voltages $V_1 = \frac{33}{7}\,\text{V}$ and $V_2 = \frac{12}{7}\,\text{V}$. The resulting voltage $v_x$ was calculated as:
        \[
        v_x = V_2 - V_{\text{bottom-right}} = \frac{12}{7} - \frac{132}{7} = -\frac{120}{7}\,\text{V}.
        \]
        This matches our result from Superposition exactly.
\end{answer}

\end{partquestion}

\end{mainquestion}
